% =========================================================
% Structures, Interpretations, and Variable Assignments
% =========================================================

\subsection{Semantics: Structures and Variable Assignments}

\begin{remark*}[Section toolkit: Semantic setup]
This section is the first point at which semantic objects enter the Predicate
Logic chapter. A language supplies the symbols; a structure supplies a domain
and interpretations for the non-logical symbols; an assignment supplies values
for variables; term evaluation turns terms into domain elements.

\begin{center}
\small
\begin{tabular}{l l l}
\toprule
\textbf{Item} & \textbf{Notation} & \textbf{Role} \\
\midrule
First-order language & \(\mathcal L\) & symbol stock with arities \\
Structure & \(\mathcal M=\langle D,I\rangle\) & semantic object for \(\mathcal L\) \\
Variable assignment & \(s:\mathsf{Var}_{\mathcal L}\to D\) & values for variables \\
Modified assignment & \(s[x\mapsto d]\) & update at one variable \\
Term evaluation & \(\llbracket t\rrbracket_{\mathcal M,s}\) & value of a term \\
\bottomrule
\end{tabular}
\end{center}
\end{remark*}

\begin{definitionbox}{Definition (First-Order Language)}
\begin{definition}[First-Order Language]
\label{def:fo-language}
A \emph{first-order language} \(\mathcal L\) consists of variables, logical
symbols, and a signature of non-logical symbols: constant symbols, function
symbols with assigned arities, and predicate symbols with assigned arities.
Equality may be included as a distinguished logical predicate.
\end{definition}
\end{definitionbox}

\begin{remark*}[Standard quantified statement]
\[
\begin{aligned}
&\text{A first-order language } \mathcal L \text{ has symbol classes}\\
&\mathsf{Var}_{\mathcal L},\quad
\mathsf{Const}_{\mathcal L},\quad
\mathsf{Func}_{\mathcal L}^{(n)},\quad
\mathsf{Pred}_{\mathcal L}^{(n)}
\quad(n\geq 1),
\end{aligned}
\]
together with the usual logical connectives, quantifiers, and punctuation.
\end{remark*}

\begin{remark*}[Predicate reading]
\(\operatorname{FirstOrderLanguage}(L)\) means that \(L\) supplies the variable
stock, logical symbols, and arity-indexed non-logical symbols used to form
first-order terms and formulas.
\end{remark*}

\begin{remark*}[Negated quantified statement]
\[
L \text{ is not equipped with the required first-order symbol classes and arities.}
\]
\end{remark*}

\begin{remark*}[Negation predicate reading]
\(\neg\operatorname{FirstOrderLanguage}(L)\) means that \(L\) lacks some
required syntactic stock or arity assignment.
\end{remark*}

\begin{remark*}[Failure modes]
A proposed language fails when its non-logical symbols are not sorted by role,
when positive-arity symbols have no arity, or when the logical grammar is not
specified.
\end{remark*}

\begin{remark*}[Failure modes]
\[
\underbrace{\mathsf{Const}_{L}}_{\text{constant symbols}}
\qquad
\underbrace{\mathsf{Func}_{L}^{(n)}}_{\text{function symbols with arity}}
\qquad
\underbrace{\mathsf{Pred}_{L}^{(n)}}_{\text{predicate symbols with arity}}.
\]
\end{remark*}

\begin{remark*}[Interpretation]
A first-order language is still syntax. It says which expressions can be
formed, but it does not assign objects, functions, or relations to the
non-logical symbols.
\end{remark*}

\begin{remark*}[Examples]
The arithmetic language with constant symbols \(0,1\), function symbols
\(S,+,\cdot\), and predicate symbol \(<\) is a first-order language once the
arities are fixed.
\end{remark*}

\begin{dependencies}
\begin{itemize}
  \item \hyperref[def:variable]{Variable}
  \item \hyperref[def:constant-symbol]{Constant Symbol}
  \item \hyperref[def:function-symbol]{Function Symbol}
  \item \hyperref[def:predicate-symbol]{Predicate Symbol}
\end{itemize}
\end{dependencies}

\begin{definitionbox}{Definition (Structure)}
\begin{definition}[Structure]
\label{def:structure}
Let \(\mathcal L\) be a first-order language. A \emph{structure} for
\(\mathcal L\) is a pair
\[
\mathcal M=\langle D,I\rangle
\]
where \(D\) is a nonempty set and \(I\) assigns each constant symbol an element
of \(D\), each \(n\)-ary function symbol a function \(D^n\to D\), and each
\(n\)-ary predicate symbol a relation on \(D^n\).
\end{definition}
\end{definitionbox}

\begin{remark*}[Standard quantified statement]
\[
\begin{aligned}
&\text{Let } \mathcal L \text{ be a first-order language.}\\
&\mathcal M=\langle D,I\rangle \text{ is a structure for } \mathcal L\\
&\Longleftrightarrow
D\neq\varnothing,\;
I(c)\in D,\;
I(f):D^n\to D,\;
I(P)\subseteq D^n
\end{aligned}
\]
for each \(c\in\mathsf{Const}_{\mathcal L}\),
\(f\in\mathsf{Func}_{\mathcal L}^{(n)}\), and
\(P\in\mathsf{Pred}_{\mathcal L}^{(n)}\).
\end{remark*}

\begin{remark*}[Predicate reading]
\(\operatorname{Structure}(M,L)\) means that \(M\) is a structure for the
first-order language \(L\).
\end{remark*}

\begin{remark*}[Negated quantified statement]
\[
D=\varnothing
\quad\text{or}\quad
I \text{ fails to assign the required object, function, or relation to some symbol.}
\]
\end{remark*}

\begin{remark*}[Negation predicate reading]
\(\neg\operatorname{Structure}(M,L)\) means that \(M\) is not a well-formed
structure for \(L\).
\end{remark*}

\begin{remark*}[Failure modes]
A proposed structure fails if the domain is empty, if a symbol is left
uninterpreted, or if the assigned object has the wrong arity or type.
\end{remark*}

\begin{remark*}[Failure modes]
\[
\underbrace{D=\varnothing}_{\text{empty domain}}
\quad\lor\quad
\underbrace{I(c)\notin D}_{\text{bad constant value}}
\quad\lor\quad
\underbrace{I(f):D^n\not\to D}_{\text{bad function arity}}
\quad\lor\quad
\underbrace{I(P)\not\subseteq D^n}_{\text{bad predicate relation}}.
\]
\end{remark*}

\begin{remark*}[Interpretation]
A structure is the minimal semantic object needed in Volume I. It gives
meaning to the non-logical symbols of a language. Systematic topics such as
substructures, expansions, reducts, and elementary classes belong to Volume
VIII Model Theory.
\end{remark*}

\begin{remark*}[Examples]
The natural numbers with their usual interpretations of \(0,S,+,\cdot,<\)
form a structure for the arithmetic language.
\end{remark*}

\begin{dependencies}
\begin{itemize}
  \item \hyperref[def:fo-language]{First-Order Language}
\end{itemize}
\end{dependencies}

\begin{definitionbox}{Definition (Variable Assignment)}
\begin{definition}[Variable Assignment]
\label{def:var-assign}
Let \(\mathcal M=\langle D,I\rangle\) be a structure for \(\mathcal L\).
A \emph{variable assignment} for \(\mathcal M\) is a function
\[
s:\mathsf{Var}_{\mathcal L}\to D.
\]
\end{definition}
\end{definitionbox}

\begin{remark*}[Standard quantified statement]
\[
s \text{ is a variable assignment for } \mathcal M
\Longleftrightarrow
\mathrm{dom}(s)=\mathsf{Var}_{\mathcal L}
\text{ and }
\mathrm{ran}(s)\subseteq D.
\]
\end{remark*}

\begin{remark*}[Predicate reading]
\(\operatorname{VariableAssignment}(s,M)\) means that \(s\) assigns a domain
element of \(M\) to each variable of the language of \(M\).
\end{remark*}

\begin{remark*}[Negated quantified statement]
\[
\mathrm{dom}(s)\neq\mathsf{Var}_{\mathcal L}
\quad\text{or}\quad
\exists x\in\mathsf{Var}_{\mathcal L}\;(s(x)\notin D).
\]
\end{remark*}

\begin{remark*}[Negation predicate reading]
\(\neg\operatorname{VariableAssignment}(s,M)\) means that \(s\) is not a
function assigning each variable a domain element of \(M\).
\end{remark*}

\begin{remark*}[Failure modes]
An assignment fails when it leaves a variable without a value or assigns a
value outside the domain.
\end{remark*}

\begin{remark*}[Failure modes]
\[
\underbrace{\mathrm{dom}(s)\neq\mathsf{Var}_{\mathcal L}}_{\text{missing variable value}}
\quad\lor\quad
\underbrace{s(x)\notin D}_{\text{value outside domain}}.
\]
\end{remark*}

\begin{remark*}[Interpretation]
Structures interpret non-logical symbols; assignments interpret variables.
Open formulas require assignments, while sentences eventually do not depend on
the particular assignment.
\end{remark*}

\begin{dependencies}
\begin{itemize}
  \item \hyperref[def:structure]{Structure}
\end{itemize}
\end{dependencies}

\begin{definitionbox}{Definition (Modified Assignment)}
\begin{definition}[Modified Assignment]
\label{def:mod-assign}
Let \(s:\mathsf{Var}_{\mathcal L}\to D\), let \(x\in\mathsf{Var}_{\mathcal L}\),
and let \(d\in D\). The \emph{modified assignment} \(s[x\mapsto d]\) is the
assignment defined by
\[
s[x\mapsto d](y)=
\begin{cases}
d, & y=x,\\
s(y), & y\neq x.
\end{cases}
\]
\end{definition}
\end{definitionbox}

\begin{remark*}[Standard quantified statement]
\[
\forall y\in\mathsf{Var}_{\mathcal L}\quad
s[x\mapsto d](y)=d \Longleftrightarrow y=x,
\qquad
s[x\mapsto d](y)=s(y) \Longleftrightarrow y\neq x.
\]
\end{remark*}

\begin{remark*}[Predicate reading]
\(\operatorname{ModifiedAssignment}(s,x,d,s')\) means that \(s'\) is the
assignment obtained from \(s\) by sending \(x\) to \(d\) and leaving every
other variable unchanged.
\end{remark*}

\begin{remark*}[Negated quantified statement]
\[
s'(x)\neq d
\quad\text{or}\quad
\exists y\in\mathsf{Var}_{\mathcal L}\;(y\neq x\land s'(y)\neq s(y)).
\]
\end{remark*}

\begin{remark*}[Negation predicate reading]
\(\neg\operatorname{ModifiedAssignment}(s,x,d,s')\) means that \(s'\) is not the
one-variable update of \(s\) at \(x\) with value \(d\).
\end{remark*}

\begin{remark*}[Failure modes]
A proposed modified assignment fails if it does not change \(x\) to \(d\), or
if it changes some variable other than \(x\).
\end{remark*}

\begin{remark*}[Failure modes]
\[
\underbrace{s'(x)\neq d}_{\text{target variable not updated}}
\quad\lor\quad
\underbrace{\exists y\neq x\;(s'(y)\neq s(y))}_{\text{unwanted extra change}}.
\]
\end{remark*}

\begin{remark*}[Interpretation]
Modified assignments are the semantic mechanism behind quantifiers. To test a
quantified formula, the assignment is varied at the bound variable while all
other variable values remain fixed.
\end{remark*}

\begin{dependencies}
\begin{itemize}
  \item \hyperref[def:var-assign]{Variable Assignment}
\end{itemize}
\end{dependencies}

\begin{definitionbox}{Definition (Term Evaluation)}
\begin{definition}[Term Evaluation]
\label{def:term-interp}
Let \(\mathcal M=\langle D,I\rangle\) be a structure and \(s\) a variable
assignment. The \emph{value} of a term \(t\), denoted
\(\llbracket t\rrbracket_{\mathcal M,s}\), is defined recursively by
\[
\llbracket x\rrbracket_{\mathcal M,s}=s(x),\qquad
\llbracket c\rrbracket_{\mathcal M,s}=I(c),
\]
and
\[
\llbracket f(t_1,\ldots,t_n)\rrbracket_{\mathcal M,s}
=
I(f)\bigl(\llbracket t_1\rrbracket_{\mathcal M,s},\ldots,
\llbracket t_n\rrbracket_{\mathcal M,s}\bigr).
\]
\end{definition}
\end{definitionbox}

\begin{remark*}[Standard quantified statement]
\[
\begin{aligned}
&\llbracket x\rrbracket_{\mathcal M,s}=s(x),\\
&\llbracket c\rrbracket_{\mathcal M,s}=I(c),\\
&\llbracket f(t_1,\ldots,t_n)\rrbracket_{\mathcal M,s}
=I(f)(\llbracket t_1\rrbracket_{\mathcal M,s},\ldots,
\llbracket t_n\rrbracket_{\mathcal M,s}).
\end{aligned}
\]
\end{remark*}

\begin{remark*}[Predicate reading]
\(\operatorname{TermEvaluation}(M,s,t,a)\) means that the term \(t\) evaluates
to the domain element \(a\) in \(M\) under \(s\).
\end{remark*}

\begin{remark*}[Negated quantified statement]
\[
\llbracket t\rrbracket_{\mathcal M,s}\neq a.
\]
\end{remark*}

\begin{remark*}[Negation predicate reading]
\(\neg\operatorname{TermEvaluation}(M,s,t,a)\) means that \(a\) is not the value
of \(t\) in \(M\) under \(s\).
\end{remark*}

\begin{remark*}[Failure modes]
Term evaluation fails to have the proposed value when a variable, constant, or
function-application clause computes a different domain element.
\end{remark*}

\begin{remark*}[Failure modes]
\[
\underbrace{s(x)\neq a}_{\text{variable case}}
\quad\lor\quad
\underbrace{I(c)\neq a}_{\text{constant case}}
\quad\lor\quad
\underbrace{I(f)(a_1,\ldots,a_n)\neq a}_{\text{function case}}.
\]
\end{remark*}

\begin{remark*}[Interpretation]
Term evaluation is the semantic counterpart of term formation. It assigns a
domain element to each grammatical term once a structure and assignment are
fixed.
\end{remark*}

\begin{dependencies}
\begin{itemize}
  \item \hyperref[def:term]{Term}
  \item \hyperref[def:structure]{Structure}
  \item \hyperref[def:var-assign]{Variable Assignment}
\end{itemize}
\end{dependencies}
